\begin{center}\itshape
Walk up a vertical line to the right of $\frac12$ and record the values of $\zeta$.\\
They wander over the whole plane, with a law that you can compute,\\
because the primes, seen from far away, are independent.
\end{center}

Few analytic number theorists own an Olympic medal. Harald Bohr does: he played for the Danish football team that won silver at the London Olympics of 1908, two years before his doctoral thesis. His thesis was about Dirichlet series, and in the years that followed, with Richard Courant and Edmund Landau, he studied the values of $\zeta(\sigma+it)$ on a vertical line $\sigma>\frac12$, which nobody before him had thought of as a single object. In 1914 Bohr and Courant proved that for $\frac12<\sigma\le1$ these values are dense in the complex plane: $\zeta$ comes as close as you like to every complex number. Figure~\ref{fig:zetacurve} shows the curve $t\mapsto\zeta(0.7+it)$ beginning its wandering.

The idea behind Bohr's results is the theme of this book in its purest form. The Euler product $\zeta(s)=\prod_p(1-p^{-s})^{-1}$ depends on $t$ only through the numbers $p^{-it}=e^{-it\log p}$, points on the unit circle. Because the numbers $\log p$ are linearly independent over $\mathbb Q$ (that is unique factorisation), these points behave, for a random $t$, like \emph{independent} uniform points: the Steinhaus random multiplicative function of Chapter~\ref{ch:randommult}. In 1930 and 1932 Bohr and Børge Jessen turned this into a limit theorem: $\zeta(\sigma+it)$, for $t$ uniform in a long interval, converges in law to a random Euler product. This chapter proves it, deduces the density of the values, and shows that the limit law of $\log|\zeta|$ has a smooth density. On the critical line itself the random Euler product diverges, and Section~\ref{ch:selbergclt} takes over.

Throughout, $\frac12<\sigma<1$ is fixed, $t$ is uniform on $[T,2T]$, and $\mathbb E_T$ and $P_T$ refer to this choice. $(X_p)_p$ are independent and uniform on the unit circle, and for $|x|=1$,
\[
\ell_p(x)=-\log\Bigl(1-\frac x{p^\sigma}\Bigr)=\sum_{k\ge1}\frac{x^k}{kp^{k\sigma}}\qquad(\text{principal branch; }|xp^{-\sigma}|\le2^{-1/2}).
\]

\begin{reminder}[from Volumes~4, 7, 11 and~14]
\emph{Stone--Weierstrass} (Volume~7, Theorem~4.19). On a compact space, an algebra of continuous complex functions that separates points, contains the constants and is stable under conjugation is dense for the uniform norm.
\emph{The three-series theorem} (Volume~11, Theorem~2.5). A series of independent, centred random variables bounded by $c$ converges almost surely if the sum of their variances is finite.
\emph{$\zeta$ as a short sum} (Volume~14, Lemma~1.5). For integers $k\ge2$, $N\ge1$ and $s=\sigma+it\ne1$, $\sigma>0$: if $N\ge(|s|+k)^{k/(k-1)}$, then $\zeta(s)=\sum_{n\le N}n^{-s}+\frac{N^{1-s}}{s-1}+E$ with $|E|\le\frac{k+1}2N^{-\sigma}$.
\emph{Inversion} (Volume~4, Theorem~6.8). A real random variable whose characteristic function $\varphi$ is integrable has the continuous density $\frac1{2\pi}\int e^{-iux}\varphi(u)\dd u$.
\emph{Riemann--Lebesgue} (Volume~7, Exercise~8.8). $\int f(x)e^{iux}\dd x\to0$ as $|u|\to\infty$ for integrable $f$.
\end{reminder}

\begin{figure}[ht]
\centering\fig{ch06-curve}
\caption{The curve $t\mapsto\zeta(0.7+it)$ for $0\le t\le60$, starting from $\zeta(0.7)=-2.78$ (rust dot). It loops around the origin without passing through it ($\zeta$ has no zeros on this line, if the Riemann hypothesis holds), and by Corollary~\ref{cor:dense} it will eventually pass close to every point of the plane.}
\label{fig:zetacurve}
\end{figure}

\section{The random Euler product}

\begin{proposition}\label{prop:randomeuler}
Almost surely the series $L(\sigma,X)=\sum_p\ell_p(X_p)$ converges, and so does the product $\zeta(\sigma,X)=\prod_p(1-X_pp^{-\sigma})^{-1}=e^{L(\sigma,X)}$, which is therefore non-zero. Moreover $\mathbb E\ell_p(X_p)=0$ and $\mathbb E|\ell_p(X_p)|^2=\sum_{k\ge1}\frac1{k^2p^{2k\sigma}}$.
\end{proposition}
\begin{proof}
Write $\ell_p(X_p)=X_pp^{-\sigma}+r_p$ with $|r_p|\le\sum_{k\ge2}\frac{p^{-k\sigma}}k\le\frac{p^{-2\sigma}}{2(1-2^{-1/2})}\le2p^{-2\sigma}$. Since $2\sigma>1$, $\sum_p|r_p|$ converges. The real parts $\Re(X_pp^{-\sigma})$ are independent, centred, bounded by $1$, with variances $\frac12p^{-2\sigma}$, whose sum is finite; so $\sum_p\Re(X_pp^{-\sigma})$ converges almost surely by the three-series theorem, and likewise the imaginary parts. Hence $L(\sigma,X)$ converges, and $\prod_{p\le y}(1-X_pp^{-\sigma})^{-1}=\exp\sum_{p\le y}\ell_p(X_p)\to e^{L(\sigma,X)}$. Finally $\mathbb EX_p^k=0$ and $\mathbb EX_p^k\overline{X_p^j}=\mathbf 1\{j=k\}$ for $j,k\ge1$, which gives the two moments (the series converges uniformly in $X_p$).
\end{proof}

\begin{lemma}[Kronecker--Weyl for the primes]\label{lem:kw}
Let $p_1,\dots,p_r$ be distinct primes. As $T\to\infty$, $(p_1^{-it},\dots,p_r^{-it})$ converges in law to $(X_{p_1},\dots,X_{p_r})$: for every continuous $F$ on the torus $\mathbb T^r=\{|z_1|=\dots=|z_r|=1\}$, $\mathbb E_TF(p_1^{-it},\dots,p_r^{-it})\to\mathbb EF(X_{p_1},\dots,X_{p_r})$.
\end{lemma}
\begin{proof}
For $m\in\Z^r$ and the monomial $F(z)=z_1^{m_1}\cdots z_r^{m_r}$, $F(p_1^{-it},\dots)=q^{-it}$ with $q=\prod_jp_j^{m_j}$, a positive rational. If $m=0$, both sides are $1$. If $m\ne0$, then $q\ne1$ by unique factorisation, so $|\mathbb E_Tq^{-it}|=|\frac1T\int_T^{2T}e^{-it\log q}\dd t|\le\frac2{T|\log q|}\to0$, while $\mathbb E\prod_jX_{p_j}^{m_j}=0$. By linearity the convergence holds for trigonometric polynomials, which form an algebra that contains the constants, is stable under conjugation ($\overline{z_j}=z_j^{-1}$ on the torus) and separates points. By Stone--Weierstrass they are uniformly dense in $C(\mathbb T^r)$, and a uniform approximation within $\varepsilon$ changes both sides by at most $\varepsilon$.
\end{proof}

\begin{think}
Where did the proof use that the $p_j$ are primes? Try $(2^{-it},4^{-it})$: what is its limit law, and why is it not that of two independent uniform points? (Exercise~\ref{ex:kwfail}.)
\end{think}

\section{The mean square of a mollified zeta}

To pass from finitely many primes to $\zeta$, we show that $\zeta(\sigma+it)$ is close, in mean square, to its Euler product over the primes up to $y$. The tool is a crude mean value theorem for Dirichlet polynomials.

\begin{lemma}\label{lem:crudemvt}
For complex $a_1,\dots,a_N$ and $T>0$,
\[
\Bigl|\int_T^{2T}\Bigl|\sum_{n\le N}a_nn^{-it}\Bigr|^2\dd t-T\sum_{n\le N}|a_n|^2\Bigr|\le4\sum_{m<n\le N}|a_ma_n|\frac n{n-m}.
\]
If moreover $|a_n|\le Kn^{-\sigma}$ with $\frac12<\sigma<1$, the right side is at most $\frac{32K^2}{(1-\sigma)(2-2\sigma)}N^{2-2\sigma}(1+\log N)$.
\end{lemma}
\begin{proof}
Expanding the square, the diagonal terms give $T\sum|a_n|^2$, and for $m<n$ the two terms $(m,n)$ and $(n,m)$ contribute at most $2|a_ma_n||\int_T^{2T}e^{it\log(n/m)}\dd t|\le\frac{4|a_ma_n|}{\log(n/m)}$, and $\log\frac nm\ge1-\frac mn=\frac{n-m}n$. For the second claim, fix $n$ and split the sum over $m<n$. For $m\le\frac n2$, $\frac n{n-m}\le2$ and $\sum_{m\le n/2}m^{-\sigma}\le\frac{n^{1-\sigma}}{1-\sigma}$, which gives at most $\frac{2n^{1-2\sigma}}{1-\sigma}$ after multiplying by $n^{-\sigma}$. For $\frac n2<m<n$, $m^{-\sigma}\le2n^{-\sigma}$ and $\sum_{n/2<m<n}\frac n{n-m}\le n(1+\log n)$, which gives at most $2n^{1-2\sigma}(1+\log n)$. With the factors $K^2$, the sum over $m<n$ is at most $\frac{4K^2n^{1-2\sigma}(1+\log n)}{1-\sigma}$. Finally $\sum_{n\le N}n^{1-2\sigma}\le1+\frac{N^{2-2\sigma}}{2-2\sigma}\le\frac{2N^{2-2\sigma}}{2-2\sigma}$, and $4\cdot4\cdot2=32$.
\end{proof}

Let $A_y(s)=\prod_{p\le y}(1-p^{-s})$, a Dirichlet polynomial: $A_y(s)=\sum_{d\mid P_y}\mu(d)d^{-s}$ with $P_y=\prod_{p\le y}p$.

\begin{proposition}\label{prop:mollified}
For $y\ge2$,
\[
\limsup_{T\to\infty}\ \mathbb E_T\bigl|\zeta(\sigma+it)A_y(\sigma+it)-1\bigr|^2\le\sum_{n>y}\frac1{n^{2\sigma}}.
\]
\end{proposition}
\begin{proof}
Choose $\eta>0$ with $(1+\eta)(2-2\sigma)<1$, which is possible because $\sigma>\frac12$, then an integer $k\ge2$ with $\frac k{k-1}\le1+\eta$, and let $N=\lceil(3T+k)^{1+\eta}\rceil$. For $t\in[T,2T]$ and $T\ge2$, $|s|+k\le3T+k$, so Volume~14, Lemma~1.5 gives $\zeta(s)=Z_N(s)+R(s)$ with $Z_N(s)=\sum_{n\le N}n^{-s}$ and
\[
|R(s)|\le\frac{N^{1-\sigma}}{T}+\frac{k+1}2N^{-\sigma}\longrightarrow0\quad\text{uniformly, since }(1+\eta)(1-\sigma)<1.
\]
Since $|A_y(s)|\le\prod_{p\le y}2=2^{\pi(y)}$, it suffices to show that $\limsup\mathbb E_T|Z_NA_y-1|^2\le\sum_{n>y}n^{-2\sigma}$ (then use $|u+v|^2\le(1+\delta)|u|^2+(1+\frac1\delta)|v|^2$ and let $\delta\to0$).

Multiplying out, $Z_N(s)A_y(s)=\sum_{n\le NP_y}c_nn^{-s}$ with $c_n=\sum\mu(d)$ over the $d\mid\gcd(n,P_y)$ with $\frac nd\le N$. For $n\le N$ the condition $\frac nd\le N$ is automatic, and $c_n=\sum_{d\mid\gcd(n,P_y)}\mu(d)$ is $1$ if $n$ has no prime factor $\le y$ and $0$ otherwise. In particular $c_1=1$, and for $1<n\le N$, $c_n\ne0$ forces $n>y$. For all $n$, $|c_n|\le2^{\pi(y)}$. So $Z_NA_y-1=\sum_{2\le n\le NP_y}a_nn^{-it}$ with $a_n=c_nn^{-\sigma}$, and by Lemma~\ref{lem:crudemvt} with $K=2^{\pi(y)}$,
\[
\mathbb E_T|Z_NA_y-1|^2\le\sum_{y<n\le N}\frac1{n^{2\sigma}}+4^{\pi(y)}\sum_{n>N}\frac1{n^{2\sigma}}+\frac{C_{\sigma}4^{\pi(y)}}T(NP_y)^{2-2\sigma}(1+\log NP_y).
\]
As $T\to\infty$ with $y$ fixed, $N\to\infty$, so the middle term tends to $0$, and the last is $O_{y,\sigma}(T^{(1+\eta)(2-2\sigma)-1}\log T)\to0$ by the choice of $\eta$.
\end{proof}

The proof shows the mechanism clearly: $\zeta A_y-1$ is a Dirichlet series over the integers with no prime factor up to $y$, and its mean square is the sum of the squares of its coefficients, as if the $n^{-it}$ were orthogonal. The only role of $\sigma>\frac12$ is to make that sum, $\sum_{n>y}n^{-2\sigma}$, finite and small.

\section{The theorem of Bohr and Jessen}

\begin{theorem}[Bohr and Jessen, 1930--1932]\label{thm:bohrjessen}
Let $\frac12<\sigma<1$. As $T\to\infty$, $\zeta(\sigma+it)$ with $t$ uniform on $[T,2T]$ converges in law to $\zeta(\sigma,X)$: for every bounded Lipschitz $g:\C\to\R$,
\[
\mathbb E_Tg\bigl(\zeta(\sigma+it)\bigr)\longrightarrow\mathbb Eg\bigl(\zeta(\sigma,X)\bigr).
\]
Consequently $\log|\zeta(\sigma+it)|$ converges in law to $\Re L(\sigma,X)=\log|\zeta(\sigma,X)|$.
\end{theorem}
\begin{proof}
Let $g$ be bounded by $1$ and $1$-Lipschitz (the general case follows by scaling). For $y\ge2$ let $E_y(t)=A_y(\sigma+it)^{-1}=\prod_{p\le y}(1-p^{-\sigma-it})^{-1}$ and $\zeta_y(X)=\prod_{p\le y}(1-X_pp^{-\sigma})^{-1}$.

\emph{Finitely many primes.} $E_y(t)=\Phi(p^{-it})_{p\le y}$ for the continuous function $\Phi(z)=\prod_{p\le y}(1-z_pp^{-\sigma})^{-1}$ on the torus, so by Lemma~\ref{lem:kw}, $\mathbb E_Tg(E_y)\to\mathbb Eg(\zeta_y(X))$ and $\mathbb E_T|E_y|^2\to\mathbb E|\zeta_y(X)|^2$. By independence and Parseval, $\mathbb E|1-X_pp^{-\sigma}|^{-2}=\mathbb E|\sum_kX_p^kp^{-k\sigma}|^2=\sum_kp^{-2k\sigma}=(1-p^{-2\sigma})^{-1}$, so $\mathbb E|\zeta_y(X)|^2=\prod_{p\le y}(1-p^{-2\sigma})^{-1}\le\zeta(2\sigma)$.

\emph{Comparison.} Since $\zeta=E_y\cdot(\zeta A_y)$, by Cauchy--Schwarz,
\[
\bigl|\mathbb E_Tg(\zeta)-\mathbb E_Tg(E_y)\bigr|\le\mathbb E_T|\zeta-E_y|=\mathbb E_T\bigl[|E_y||\zeta A_y-1|\bigr]\le\bigl(\mathbb E_T|E_y|^2\bigr)^{1/2}\bigl(\mathbb E_T|\zeta A_y-1|^2\bigr)^{1/2}.
\]
By the first step and Proposition~\ref{prop:mollified},
\[
\limsup_{T\to\infty}\bigl|\mathbb E_Tg(\zeta)-\mathbb Eg(\zeta_y(X))\bigr|\le\zeta(2\sigma)^{1/2}\varepsilon_y,\qquad\varepsilon_y=\Bigl(\sum_{n>y}n^{-2\sigma}\Bigr)^{1/2}.
\]

\emph{All primes.} By Proposition~\ref{prop:randomeuler}, $\zeta_y(X)\to\zeta(\sigma,X)$ almost surely, so $\mathbb Eg(\zeta_y(X))\to\mathbb Eg(\zeta(\sigma,X))$ by dominated convergence. Letting $y\to\infty$, $\limsup_T|\mathbb E_Tg(\zeta)-\mathbb Eg(\zeta(\sigma,X))|\le0$.

\emph{Logarithms.} The function $z\mapsto\log|z|$ is continuous on $\C\setminus\{0\}$ and $\zeta(\sigma,X)\ne0$ almost surely. For bounded Lipschitz $h:\R\to\R$ and $\delta>0$, let $g_\delta(z)=h(\log\max(|z|,\delta))$, which is bounded and Lipschitz on $\C$. Then $|\mathbb E_Th(\log|\zeta|)-\mathbb E_Tg_\delta(\zeta)|\le2\|h\|_\infty P_T(|\zeta|<\delta)$ and similarly for $\zeta(\sigma,X)$. Since $P_T(|\zeta|<\delta)\le\mathbb E_T\phi_\delta(\zeta)\to\mathbb E\phi_\delta(\zeta(\sigma,X))\le P(|\zeta(\sigma,X)|<2\delta)$ for the Lipschitz function $\phi_\delta(z)=\min(1,(2-\frac{|z|}\delta)^+)$, and $P(|\zeta(\sigma,X)|<2\delta)\to0$ as $\delta\to0$, the claim follows.
\end{proof}

A remarkable feature of the proof is that it never mentions the zeros of $\zeta$: if $\zeta$ had zeros with real part $\sigma$, they would occupy a set of $t$ of measure tending to $0$ in proportion (Volume~14, Corollary~1.8 already says they are rare), and the theorem would not see them. Figure~\ref{fig:bohrjessen} compares the two laws numerically.

\begin{figure}[ht]
\centering\fig{ch06-density}
\caption{Histogram of $\log|\zeta(0.7+it)|$ at $10\,000$ random points $t\in[10^5,2\cdot10^5]$ (bars), against the density of $\log|\zeta(0.7,X)|$ estimated from $40\,000$ samples of the random Euler product over $p\le10^5$ (rust). The Kolmogorov--Smirnov distance between the two samples is $0.0095$, within sampling error; the variances are $0.525$ and $0.531$, and the exact variance $\frac12\sum_p\sum_k\frac1{k^2p^{1.4k}}$ is $0.529$ (Exercise~\ref{ex:varbj}).}
\label{fig:bohrjessen}
\end{figure}

\section{Every value, and a smooth density}

\begin{lemma}\label{lem:support}
For every $w\in\C$ and $\varepsilon>0$, $P(|L(\sigma,X)-w|<\varepsilon)>0$.
\end{lemma}
\begin{proof}
Write $\ell_p(x)=xp^{-\sigma}+r_p(x)$ with $|r_p(x)|\le2p^{-2\sigma}$, as in Proposition~\ref{prop:randomeuler}. Choose $y_0$ such that, for $p>y_0$, $p^{-\sigma}<\frac\varepsilon4$, and such that $\sum_{p>y_0}2p^{-2\sigma}<\frac\varepsilon4$ and $\sum_{p>y_0}\mathbb E|\ell_p(X_p)|^2<\frac{\varepsilon^2}{32}$.

\emph{A deterministic choice.} Let $x_p=1$ for $p\le y_0$, and $w'=w-\sum_{p\le y_0}\ell_p(1)$. For the primes $p>y_0$ in increasing order, let $D$ be the current value of $w'-\sum_{y_0<q<p}\ell_q(x_q)$ and choose $x_p=\frac D{|D|}$ (or $x_p=1$ if $D=0$). If $|D|\ge p^{-\sigma}$, the new value has modulus at most $|D|-p^{-\sigma}+2p^{-2\sigma}$. Since $\sum_pp^{-\sigma}=\infty$ for $\sigma<1$ and $\sum_pp^{-2\sigma}<\infty$, this cannot happen for every $p$; at the first prime $p=y_1$ with $|D|<p^{-\sigma}$, the new value has modulus at most $p^{-\sigma}+2p^{-2\sigma}<\frac\varepsilon2$. So $|\sum_{p\le y_1}\ell_p(x_p)-w|<\frac\varepsilon2$.

\emph{Positive probability.} The map $(z_p)_{p\le y_1}\mapsto\sum_{p\le y_1}\ell_p(z_p)$ is continuous on the torus, so there is $\delta>0$ such that $|\sum_{p\le y_1}\ell_p(X_p)-w|<\frac{3\varepsilon}4$ whenever $|X_p-x_p|<\delta$ for all $p\le y_1$. This event $B$ has probability $\prod_{p\le y_1}P(|X_p-x_p|<\delta)>0$. The tail $R=\sum_{p>y_1}\ell_p(X_p)$ is independent of $B$, centred, with $\mathbb E|R|^2=\sum_{p>y_1}\mathbb E|\ell_p(X_p)|^2<\frac{\varepsilon^2}{32}$ (independent centred terms), so $P(|R|\ge\frac\varepsilon4)\le\frac12$ by Chebyshev. Hence $P(|L(\sigma,X)-w|<\varepsilon)\ge\frac12P(B)>0$.
\end{proof}

\begin{corollary}[Bohr and Courant, 1914]\label{cor:dense}
Let $\frac12<\sigma<1$. For every $z\in\C$ and $\varepsilon>0$ there are arbitrarily large $t$ with $|\zeta(\sigma+it)-z|<\varepsilon$. In fact the set of such $t$ in $[T,2T]$ has measure at least $cT$ for some $c=c(z,\varepsilon,\sigma)>0$ and all large $T$.
\end{corollary}
\begin{proof}
We may assume $z\ne0$ (approach $0$ through non-zero points) and $\varepsilon<|z|$. Write $z=e^w$. By continuity of $\exp$ at $w$ there is $\varepsilon'>0$ with $|e^v-z|<\frac\varepsilon2$ when $|v-w|<\varepsilon'$, so $q=P(|\zeta(\sigma,X)-z|<\frac\varepsilon2)>0$ by Lemma~\ref{lem:support}. Let $g(u)=\min(1,(2-\frac{2|u-z|}\varepsilon)^+)$, which is Lipschitz, equals $1$ on $|u-z|\le\frac\varepsilon2$ and vanishes off $|u-z|<\varepsilon$. By Theorem~\ref{thm:bohrjessen}, $P_T(|\zeta(\sigma+it)-z|<\varepsilon)\ge\mathbb E_Tg(\zeta)\to\mathbb Eg(\zeta(\sigma,X))\ge q$, so for large $T$ the measure of the set is at least $\frac q2T$.
\end{proof}

\begin{beauty}
Bohr's theorem is about a single deterministic function, and its proof is a statement about probability: the values of $\zeta$ come close to $z$ because a random Euler product does so with positive probability, and the primes, through Kronecker--Weyl, are random enough. No explicit $t$ is ever found. In 1975 Sergei Voronin went much further: every non-vanishing holomorphic function on a small disc inside the strip is approximated by a vertical translate $\zeta(s+it)$ of zeta. Bhaskar Bagchi showed in 1981 that this \emph{universality} is again a support theorem for a random Euler product, now viewed as a random holomorphic function.
\end{beauty}

\begin{theorem}[Jessen and Wintner, 1935]\label{thm:smoothdensity}
The random variable $\log|\zeta(\sigma,X)|$ has a density of class $C^\infty$ on $\R$. Its characteristic function $\varphi(u)$ satisfies $|\varphi(u)|\le\exp(-c|u|^{1/\sigma}/\log|u|)$ for $|u|\ge u_0$, with $c,u_0>0$ depending on $\sigma$.
\end{theorem}
\begin{proof}
\emph{Factorisation.} By independence and the almost sure convergence of $\sum_p\Re\ell_p(X_p)$, $\varphi(u)=\mathbb Ee^{iu\log|\zeta(\sigma,X)|}=\prod_p\varphi_p(u)$ with $\varphi_p(u)=\mathbb Ee^{iu\Re\ell_p(X_p)}$, and $|\varphi_p(u)|\le1$ (dominated convergence for the partial products).

\emph{One prime.} Let $J(v)=\mathbb Ee^{iv\Re X_p}=\frac1{2\pi}\int_0^{2\pi}e^{iv\cos\theta}\dd\theta=\int_{-1}^1e^{ivx}\frac{\dd x}{\pi\sqrt{1-x^2}}$ (the real part of a uniform point of the circle has the arcsine density on $[-1,1]$). By Riemann--Lebesgue, $J(v)\to0$; fix $A\ge1$ with $|J(v)|\le\frac14$ for $v\ge A$. Since $|e^{ia}-e^{ib}|\le|a-b|$ and $|\Re\ell_p(X_p)-\Re X_pp^{-\sigma}|\le2p^{-2\sigma}$,
\[
|\varphi_p(u)|\le|J(|u|p^{-\sigma})|+2|u|p^{-2\sigma}.
\]

\emph{Many primes.} Let $u\ge u_0$ (the case $u<0$ follows from $\varphi(-u)=\overline{\varphi(u)}$), and consider the primes with $A\le up^{-\sigma}\le2A$, that is $(\frac u{2A})^{1/\sigma}\le p\le(\frac uA)^{1/\sigma}$. For these, $2up^{-2\sigma}\le4Ap^{-\sigma}\le\frac14$ once $u$ is large, so $|\varphi_p(u)|\le\frac12$. By the prime number theorem their number is at least $c_1\frac{u^{1/\sigma}}{\log u}$ for $u\ge u_0$, with $c_1=c_1(A,\sigma)>0$, because the interval is $[aU,bU]$ with $U=u^{1/\sigma}$ and fixed $0<a<b$. Bounding the other factors by $1$, $|\varphi(u)|\le2^{-c_1u^{1/\sigma}/\log u}$.

\emph{Smoothness.} Hence $\int|u|^k|\varphi(u)|\dd u<\infty$ for every $k$. By Volume~4, Theorem~6.8, the law has the continuous density $f(v)=\frac1{2\pi}\int e^{-iuv}\varphi(u)\dd u$, and differentiating $k$ times under the integral sign (justified by dominated convergence with the integrable bound $|u|^k|\varphi(u)|$) shows that $f$ is $C^k$ for every $k$.
\end{proof}

\begin{remark}
These are quoted for culture; see the Sources. (a) Jessen and Wintner (1935) proved that the law of the complex variable $L(\sigma,X)$ also has a smooth density in the plane. (b) The upper tail is much thinner than Gaussian: Lamzouri proved that the proportion of $t\in[T,2T]$ with $\log|\zeta(\sigma+it)|\ge\tau$ is $\exp\bigl(-A(\sigma)\tau^{1/(1-\sigma)}(\log\tau)^{\sigma/(1-\sigma)}(1+o(1))\bigr)$, uniformly for $\tau$ up to a constant times $(\log T)^{1-\sigma}/\log\log T$, the same tail as for the random Euler product. Montgomery (1977) had shown that values this large do occur. (c) Voronin's universality theorem (1975) and Bagchi's thesis (1981) described in the box above.
\end{remark}
